% !TEX program = pdflatex
\documentclass[11pt,a4paper]{article}

% ── Encoding & Language ──
\usepackage[utf8]{inputenc}
\usepackage[T1]{fontenc}
\usepackage[italian]{babel}

% ── Page geometry ──
\usepackage[margin=2.2cm,top=2cm,bottom=2cm]{geometry}

% ── Typography ──
\usepackage{lmodern}
\usepackage{microtype}

% ── Colours ──
\usepackage[dvipsnames,table]{xcolor}
\definecolor{codebg}{HTML}{F5F5F5}
\definecolor{codeframe}{HTML}{CCCCCC}
\definecolor{linkblue}{HTML}{0366D6}
\definecolor{notebg}{HTML}{FFF8DC}
\definecolor{noteframe}{HTML}{DAA520}
\definecolor{boxblue}{HTML}{E8F0FE}
\definecolor{boxblueframe}{HTML}{4285F4}
\definecolor{ideagreen}{HTML}{E6F4EA}
\definecolor{ideagreenframe}{HTML}{34A853}

% ── Graphics, tables, lists ──
\usepackage{graphicx}
\usepackage{booktabs}
\usepackage{longtable}
\usepackage{enumitem}
\usepackage{multicol}
\usepackage{float}

% ── Math ──
\usepackage{amsmath,amssymb}

% ── Code listings ──
\usepackage{listings}
\lstset{
  backgroundcolor=\color{codebg},
  frame=single,
  rulecolor=\color{codeframe},
  basicstyle=\ttfamily\small,
  keywordstyle=\bfseries\color{MidnightBlue},
  commentstyle=\itshape\color{ForestGreen},
  stringstyle=\color{BrickRed},
  numberstyle=\tiny\color{gray},
  numbers=left,
  numbersep=6pt,
  breaklines=true,
  breakatwhitespace=false,
  tabsize=4,
  showstringspaces=false,
  captionpos=b,
  xleftmargin=1em,
  framexleftmargin=0.5em,
  aboveskip=0.8em,
  belowskip=0.6em,
}
\lstdefinestyle{cpp}{language=C++,morekeywords={uint64_t,uint32_t,size_t,auto,std,constexpr,alignas,inline}}
\lstdefinestyle{bash}{language=bash,morekeywords={sbatch,srun,ssh,nproc}}
\lstdefinestyle{make}{language=make}
\lstdefinestyle{plain}{numbers=none,frame=none,backgroundcolor=\color{white},basicstyle=\ttfamily\small}

% ── Boxes (tcolorbox) ──
\usepackage[most]{tcolorbox}
\newtcolorbox{notebox}[1][]{
  colback=notebg, colframe=noteframe, fonttitle=\bfseries,
  title={#1}, boxrule=0.6pt, arc=2pt, left=6pt, right=6pt,
  top=4pt, bottom=4pt
}
\newtcolorbox{ideabox}[1][]{
  colback=ideagreen, colframe=ideagreenframe, fonttitle=\bfseries,
  title={#1}, boxrule=0.6pt, arc=2pt, left=6pt, right=6pt,
  top=4pt, bottom=4pt
}
\newtcolorbox{bluebox}[1][]{
  colback=boxblue, colframe=boxblueframe, fonttitle=\bfseries,
  title={#1}, boxrule=0.6pt, arc=2pt, left=6pt, right=6pt,
  top=4pt, bottom=4pt
}

% ── Hyperlinks (load last) ──
\usepackage{hyperref}
\hypersetup{
  colorlinks=true,
  linkcolor=linkblue,
  urlcolor=linkblue,
  citecolor=linkblue,
  pdfauthor={Guida SPM Modulo 2},
  pdftitle={Guida di Sviluppo -- SPM Modulo 2},
}

% ── Header/Footer ──
\usepackage{fancyhdr}
\pagestyle{fancy}
\fancyhf{}
\fancyhead[L]{\small\textit{SPM Modulo 2 -- Guida di Sviluppo}}
\fancyhead[R]{\small\thepage}
\renewcommand{\headrulewidth}{0.4pt}

% ── Section formatting ──
\usepackage{titlesec}
\titleformat{\section}{\Large\bfseries\color{MidnightBlue}}{\thesection}{1em}{}
\titleformat{\subsection}{\large\bfseries\color{RoyalBlue}}{\thesubsection}{0.8em}{}
\titleformat{\subsubsection}{\normalsize\bfseries\color{NavyBlue}}{\thesubsubsection}{0.6em}{}

% ══════════════════════════════════════════════════════════════
\begin{document}

% ── Title page ──
\begin{titlepage}
\centering
\vspace*{3cm}
{\Huge\bfseries\color{MidnightBlue} Guida di Sviluppo\par}
\vspace{0.6cm}
{\LARGE SPM --- Modulo 2\par}
\vspace{0.4cm}
{\Large\itshape Partitioned Hash Join with Duplicates:\\
From Partition Mapping to Full Join\par}
\vspace{2cm}
\begin{tcolorbox}[colback=boxblue,colframe=boxblueframe,width=12cm,arc=4pt]
\centering
\textbf{Obiettivo}: Implementare una versione parallela (C++ threads) del Partitioned Hash Join con duplicati, partendo dal codice sequenziale fornito (\texttt{hashjoin\_seq.cpp}), e produrre un report con analisi delle performance.\\[6pt]
\textbf{Deadline}: 20 Aprile (soft deadline)\\[3pt]
\textbf{Deliverables}: codice sorgente + Makefile + README + report PDF (max 5 pagine)
\end{tcolorbox}
\vspace{1.5cm}
{\large Corso: Sistemi Paralleli e Multi-core (SPM)\\
A.Y.\ 2025--2026 --- Prof.\ Massimo Torquati}
\vfill
\end{titlepage}

% ── Table of contents ──
\tableofcontents
\clearpage

% ══════════════════════════════════════════════════════════════
\section{Panoramica dell'Algoritmo}
% ══════════════════════════════════════════════════════════════

\subsection{Il problema: cos'è un Join e perché è costoso?}

Immagina di avere due tabelle di un database:

\begin{center}
\begin{minipage}[t]{0.35\textwidth}
\centering
\textbf{Relazione R} (es.\ ``Ordini'')\\[4pt]
\begin{tabular}{|c|}
\hline
\textbf{key}\\\hline
3\\7\\3\\5\\2\\\hline
\end{tabular}
\end{minipage}
\hfill
\begin{minipage}[t]{0.35\textwidth}
\centering
\textbf{Relazione S} (es.\ ``Prodotti'')\\[4pt]
\begin{tabular}{|c|}
\hline
\textbf{key}\\\hline
5\\3\\3\\7\\1\\3\\2\\5\\\hline
\end{tabular}
\end{minipage}
\end{center}

\medskip
L'operazione di \textbf{join} trova tutte le coppie $(r,s)$ dove $r.\mathit{key} = s.\mathit{key}$. In SQL: \texttt{SELECT * FROM R JOIN S ON R.key = S.key}.

\paragraph{Approccio naive} Confronta ogni record di R con ogni record di S:
\[
\text{Complessità: } O(N_R \times N_S)
\]
Con 10 milioni di record per tabella $\Rightarrow 10^{14}$ confronti. \textbf{Inaccettabile}.

\paragraph{Hash join classico} Costruisci una hash table su R, poi probi con S:
\begin{enumerate}[nosep]
  \item Costruisci \texttt{hash\_table} da R \hfill $O(N_R)$
  \item Per ogni $s \in S$: cerca $s.\mathit{key}$ nella hash table \hfill $O(N_S)$ in media
\end{enumerate}
Complessità totale: $O(N_R + N_S)$. Ma se R ed S sono enormi, la hash table non entra in cache e ogni lookup produce \emph{cache miss} $\Rightarrow$ lento nella pratica.


\subsection{L'idea del Partitioned Hash Join: ``dividi e conquista''}

\begin{ideabox}[Idea Centrale]
Se applico la \textbf{stessa} funzione hash $h(\mathit{key})$ sia a R che a S, allora tutti i record con la stessa chiave finiranno nella \textbf{stessa partizione}.

Questo significa che posso fare il join \textbf{separatamente} su ogni partizione, perché un record in $R_{\text{partizione}_0}$ non potrà \emph{mai} matchare un record in $S_{\text{partizione}_1}$ (hanno hash diverse!).
\end{ideabox}

\medskip
Concetto visivo:

\begin{center}
\begin{minipage}[t]{0.42\textwidth}
\textbf{PRIMA} (dati mescolati):
\begin{lstlisting}[style=plain]
R = [3, 7, 3, 5, 2]
S = [5, 3, 3, 7, 1, 3, 2, 5]
\end{lstlisting}
\end{minipage}
\hfill
\begin{minipage}[t]{0.52\textwidth}
\textbf{DOPO} (partizionati con $P=3$):
\begin{lstlisting}[style=plain]
R_0 = [2]       S_0 = [1, 2]
R_1 = [3, 3, 5] S_1 = [5, 3, 3, 3, 5]
R_2 = [7]       S_2 = [7]
\end{lstlisting}
Ogni coppia $(R_p, S_p)$ è \textbf{indipendente}!
\end{minipage}
\end{center}


\subsection{Collegamento con il Modulo 1: la funzione di mapping}

Nel Modulo 1 hai implementato e ottimizzato una funzione:
\[
h(\mathit{key}) \;\rightarrow\; \mathit{partition\_id} \;\in\; \{0, 1, \dots, P{-}1\}
\]

La tua implementazione usa \textbf{XOR-fold + Fibonacci multiply-shift a 32 bit}:

\begin{lstlisting}[style=cpp,caption={Hash function dal Modulo 1}]
static constexpr uint32_t HASH_A32 = 0x9E3779B9u; // floor(2^32 / phi)

inline uint32_t compute_partition_id(uint64_t key, unsigned shift) {
    uint32_t k_lo = static_cast<uint32_t>(key);          // meta' bassa
    uint32_t k_hi = static_cast<uint32_t>(key >> 32);    // meta' alta
    uint32_t mixed = k_lo ^ k_hi;                        // XOR-fold
    return (mixed * HASH_A32) >> shift;   // shift = 32 - log2(P)
}
\end{lstlisting}

\begin{notebox}[Connessione Modulo 1 $\to$ Modulo 2]
Nel Modulo 1 hai costruito e benchmarkato questa funzione \emph{da sola} (throughput in Mkeys/s, distribuzione tra partizioni). Nel Modulo 2, questa funzione diventa il ``mattoncino'' usato \textbf{dentro} un algoritmo più grande. Viene chiamata $2 \times (N_R + N_S)$ volte (una volta nell'histogram e una nello scatter, per ciascuna relazione).

Il codice fornito dal professore usa una hash semplicistica (\lstinline{key & (P-1)}). Tu devi \textbf{sostituirla} con la tua.
\end{notebox}


\subsection{Le fasi dell'algoritmo: visione d'insieme}

\begin{figure}[H]
\centering
\begin{tcolorbox}[colback=white,colframe=MidnightBlue,title={\textbf{Partitioned Hash Join --- Pipeline}},width=\textwidth,arc=3pt]

\textbf{Macro-Fase A: Partizionamento} (applicata sia a R che a S)

\smallskip
\begin{tabular}{@{}l@{\;\;}c@{\;\;}l@{\;\;}c@{\;\;}l@{}}
\fbox{1.\ Histogram (conteggi)} & $\longrightarrow$ &
\fbox{2.\ Prefix Sum (offset)} & $\longrightarrow$ &
\fbox{3.\ Scatter (riordino)}
\end{tabular}

\medskip
\hrule
\medskip

\textbf{Macro-Fase B: Join locale} (per ogni partizione $p = 0 \dots P{-}1$, indipendente)

\smallskip
\begin{tabular}{@{}l@{\;\;}c@{\;\;}l@{}}
\fbox{4.\ Build (hash table su $R_p$)} & $\longrightarrow$ &
\fbox{5.\ Probe (cerca match da $S_p$)}
\end{tabular}

\medskip
\hrule
\medskip

\textbf{Macro-Fase C: Accumulazione} --- Somma join\_count e checksum di tutte le partizioni
\end{tcolorbox}
\end{figure}


\subsection{Esempio numerico completo (passo per passo)}
\label{sec:esempio}

Usiamo dati piccoli per tracciare l'intero algoritmo a mano.

\begin{bluebox}[Parametri dell'esempio]
$P = 4$ partizioni, \quad $h(\mathit{key}) = \mathit{key} \bmod 4$ \quad (semplificata per leggibilità)

\medskip
\begin{tabular}{l|cccccccc}
\textbf{R} & \multicolumn{8}{c}{indice: 0~~1~~2~~3~~4~~5~~6~~7}\\\hline
chiave & 5 & 2 & 8 & 3 & 2 & 5 & 7 & 4
\end{tabular}
\qquad
\begin{tabular}{l|cccccc}
\textbf{S} & \multicolumn{6}{c}{indice: 0~~1~~2~~3~~4~~5}\\\hline
chiave & 3 & 5 & 2 & 8 & 5 & 7
\end{tabular}
\end{bluebox}

\subsubsection{Step 1 --- Mapping (key $\to$ partition\_id)}

Applichiamo $h(\mathit{key}) = \mathit{key} \bmod 4$:

\begin{center}
\small
\begin{tabular}{l|cccccccc}
\textbf{R} chiave & 5 & 2 & 8 & 3 & 2 & 5 & 7 & 4\\\hline
partition\_id & 1 & 2 & 0 & 3 & 2 & 1 & 3 & 0\\
 & $5\%4$ & $2\%4$ & $8\%4$ & $3\%4$ & $2\%4$ & $5\%4$ & $7\%4$ & $4\%4$
\end{tabular}
\qquad
\begin{tabular}{l|cccccc}
\textbf{S} chiave & 3 & 5 & 2 & 8 & 5 & 7\\\hline
partition\_id & 3 & 1 & 2 & 0 & 1 & 3
\end{tabular}
\end{center}

Il mapping \textbf{non sposta nulla}: associa solo un partition\_id ad ogni record.

\subsubsection{Step 2 --- Histogram}

Scorriamo i partition\_id e contiamo quanti record per partizione:

\begin{center}
\begin{tabular}{c|cccc}
& \textbf{part 0} & \textbf{part 1} & \textbf{part 2} & \textbf{part 3}\\\hline
\texttt{hist\_R} & 2 (key 8,4) & 2 (key 5,5) & 2 (key 2,2) & 2 (key 3,7)\\
\texttt{hist\_S} & 1 (key 8) & 2 (key 5,5) & 1 (key 2) & 2 (key 3,7)
\end{tabular}
\end{center}

\textbf{A cosa serve?} Sapere quanti record ha ogni partizione permette di pre-allocare lo spazio nell'array di output e calcolare dove inizia ciascuna partizione.

\subsubsection{Step 3 --- Prefix Sum (offset di inizio)}

L'exclusive prefix sum converte i conteggi in posizioni di inizio:

\begin{center}
\begin{tabular}{l|cccc|l}
& part 0 & part 1 & part 2 & part 3 & formula\\\hline
\texttt{hist\_R}  & 2 & 2 & 2 & 2 &\\
\texttt{begin\_R} & 0 & 2 & 4 & 6 & $\text{begin}[i] = \sum_{j<i}\text{hist}[j]$\\
\texttt{end\_R}   & 2 & 4 & 6 & 8 & $\text{end}[i] = \text{begin}[i] + \text{hist}[i]$
\end{tabular}
\end{center}

Mappa delle posizioni nell'array output di R:
\begin{center}
\small
\begin{tabular}{|c|c||c|c||c|c||c|c|}
\hline
\multicolumn{2}{|c||}{\cellcolor{blue!10}part 0} &
\multicolumn{2}{c||}{\cellcolor{green!10}part 1} &
\multicolumn{2}{c||}{\cellcolor{orange!10}part 2} &
\multicolumn{2}{c|}{\cellcolor{red!10}part 3} \\\hline
pos 0 & pos 1 & pos 2 & pos 3 & pos 4 & pos 5 & pos 6 & pos 7\\\hline
\end{tabular}
\end{center}

Analogamente per S: \texttt{begin\_S} = [0, 1, 3, 4], \;\texttt{end\_S} = [1, 3, 4, 6].

\subsubsection{Step 4 --- Scatter (riordino per partizione contigua)}

Scorriamo R originale e copiamo ogni record nella posizione corretta, usando un \textbf{cursore} per partizione:

\begin{center}\small
\begin{tabular}{ll}
\toprule
\textbf{Record} & \textbf{Azione}\\\midrule
R[0] key=5, pid=1 & $\texttt{out}[\texttt{next}[1]{=}2] \leftarrow 5$, \;\texttt{next}[1]++ $\to 3$\\
R[1] key=2, pid=2 & $\texttt{out}[\texttt{next}[2]{=}4] \leftarrow 2$, \;\texttt{next}[2]++ $\to 5$\\
R[2] key=8, pid=0 & $\texttt{out}[\texttt{next}[0]{=}0] \leftarrow 8$, \;\texttt{next}[0]++ $\to 1$\\
R[3] key=3, pid=3 & $\texttt{out}[\texttt{next}[3]{=}6] \leftarrow 3$, \;\texttt{next}[3]++ $\to 7$\\
R[4] key=2, pid=2 & $\texttt{out}[\texttt{next}[2]{=}5] \leftarrow 2$, \;\texttt{next}[2]++ $\to 6$\\
R[5] key=5, pid=1 & $\texttt{out}[\texttt{next}[1]{=}3] \leftarrow 5$, \;\texttt{next}[1]++ $\to 4$\\
R[6] key=7, pid=3 & $\texttt{out}[\texttt{next}[3]{=}7] \leftarrow 7$, \;\texttt{next}[3]++ $\to 8$\\
R[7] key=4, pid=0 & $\texttt{out}[\texttt{next}[0]{=}1] \leftarrow 4$, \;\texttt{next}[0]++ $\to 2$\\
\bottomrule
\end{tabular}
\end{center}

\medskip
\textbf{Array R riordinato}:
\begin{center}\small
\begin{tabular}{|c|c||c|c||c|c||c|c|}
\hline
\multicolumn{2}{|c||}{\cellcolor{blue!10}part 0} &
\multicolumn{2}{c||}{\cellcolor{green!10}part 1} &
\multicolumn{2}{c||}{\cellcolor{orange!10}part 2} &
\multicolumn{2}{c|}{\cellcolor{red!10}part 3} \\\hline
8 & 4 & 5 & 5 & 2 & 2 & 3 & 7\\\hline
\end{tabular}
\end{center}

Analogamente, \textbf{S riordinato}: \;
\begin{tabular}{|c||c|c||c||c|c|}
\hline
\cellcolor{blue!10}8 & \cellcolor{green!10}5 & \cellcolor{green!10}5 & \cellcolor{orange!10}2 & \cellcolor{red!10}3 & \cellcolor{red!10}7\\\hline
\end{tabular}

\subsubsection{Step 5 --- Build (hash table locale su ogni $R_p$)}

Per ogni partizione scansioniamo $R_p$ e contiamo le occorrenze di ogni chiave:

\begin{center}\small
\begin{tabular}{lll}
\toprule
\textbf{Partizione} & $R_p$ & \textbf{countR}\\\midrule
$p=0$ & [8, 4] & \{8:1, 4:1\}\\
$p=1$ & [5, 5] & \{5:\textbf{2}\} $\leftarrow$ \emph{chiave 5 duplicata!}\\
$p=2$ & [2, 2] & \{2:\textbf{2}\} $\leftarrow$ \emph{chiave 2 duplicata!}\\
$p=3$ & [3, 7] & \{3:1, 7:1\}\\
\bottomrule
\end{tabular}
\end{center}

\subsubsection{Step 6 --- Probe (cercare match da $S_p$)}

Per ogni partizione scansioniamo $S_p$ e per ogni chiave cerchiamo in countR:

\begin{center}\small
\begin{tabular}{llll}
\toprule
\textbf{Part.} & $S_p$ & \textbf{Lookup} & \textbf{join\_count}\\\midrule
$p=0$ & [8] & $8 \in$ countR, mult=1 & $+1$\\
\rowcolor{green!8}
$p=1$ & [5, 5] & $5 \in$ countR, mult=2 $\times 2$ record & $+2+2 = +4$\\
$p=2$ & [2] & $2 \in$ countR, mult=2 & $+2$\\
$p=3$ & [3, 7] & $3 \in$ countR, mult=1; $7 \in$ countR, mult=1 & $+1+1 = +2$\\
\bottomrule
\end{tabular}
\end{center}

\begin{notebox}[Spiegazione della partizione 1]
La chiave 5 appare 2 volte in $R_1$ e 2 volte in $S_1$. Ogni record di S con chiave 5 contribuisce $\text{countR}[5] = 2$ match. Ci sono 2 record in $S_1$ con chiave 5, quindi il totale è $2 \times 2 = 4$ match. Questo è il meccanismo con cui l'algoritmo gestisce i \textbf{duplicati}.
\end{notebox}

\subsubsection{Step 7 --- Accumulazione}

\[
\text{join\_count} = 1 + 4 + 2 + 2 = \boxed{9}
\]

\paragraph{Verifica naive.} Elenchiamo tutte le coppie $(r,s)$ con $r.\mathit{key} = s.\mathit{key}$:

\begin{center}\small
\begin{tabular}{ccc}
R[0]=5 $\leftrightarrow$ S[1]=5 & R[3]=3 $\leftrightarrow$ S[0]=3 & R[6]=7 $\leftrightarrow$ S[5]=7 \\
R[0]=5 $\leftrightarrow$ S[4]=5 & R[4]=2 $\leftrightarrow$ S[2]=2 & \\
R[1]=2 $\leftrightarrow$ S[2]=2 & R[5]=5 $\leftrightarrow$ S[1]=5 & \\
R[2]=8 $\leftrightarrow$ S[3]=8 & R[5]=5 $\leftrightarrow$ S[4]=5 & \\
\end{tabular}
\end{center}
Totale = 9 \checkmark \; (corrisponde!)


\subsection{Perché partizionare è vantaggioso}

\begin{enumerate}[leftmargin=*]
\item \textbf{Località di cache.}
Senza partizionamento la hash table di R ha $N_R$ entry $\Rightarrow$ enorme, non sta in cache; ogni lookup $\to$ cache miss.
Con $P$ partizioni, ogni hash table locale ha $\sim N_R/P$ entry; scegliendo $P$ grande abbastanza, sta in L1/L2 $\Rightarrow$ ogni lookup è un cache hit.

\item \textbf{Indipendenza tra partizioni.}
Le $P$ partizioni non condividono dati $\Rightarrow$ nessun lock/sincronizzazione $\Rightarrow$ parallelismo naturale (\emph{embarrassingly parallel}).

\item \textbf{Gestione della memoria.}
Tante hash table piccole e di dimensione prevedibile invece di una gigante con possibili resize.
\end{enumerate}

Il costo del partizionamento è $O(N)$ per relazione; il beneficio di cache supera di gran lunga questo costo per dataset grandi.


\subsection{La gestione dei duplicati: perché ``with Duplicates''}

Se una chiave $k$ appare $m$ volte in $R$ e $n$ volte in $S$, produce $m \times n$ match:

\begin{itemize}[nosep]
  \item \textbf{Build}: conta le occorrenze $\Rightarrow$ \lstinline{countR[key]++}
  \item \textbf{Probe}: quando trova un match, aggiunge \lstinline{countR[key]} (non 1)
\end{itemize}

Esempio: key=5, $m=2$ in R, $n=2$ in S $\Rightarrow$ 4 coppie matchanti, e l'algoritmo conta correttamente $2+2=4$.


\subsection{Checksum: verifica senza materializzare le coppie}

L'algoritmo non produce la lista delle coppie (sarebbe enorme). Verifica la correttezza con due checksum deterministici:

\begin{lstlisting}[style=cpp,numbers=none]
result.checksum1 += splitmix64(key) * multiplicity;
result.checksum2 += splitmix64(key ^ 0x9e3779b97f4a7c15ULL) * multiplicity;
\end{lstlisting}

\begin{itemize}[nosep]
  \item \texttt{splitmix64} è una funzione di mixing deterministica (stesso input $\to$ stesso output)
  \item Due checksum indipendenti (salt diversi) riducono la probabilità di collisioni accidentali
  \item Se la versione parallela produce gli stessi \texttt{join\_count}, \texttt{checksum1}, \texttt{checksum2} $\Rightarrow$ correttezza verificata
\end{itemize}


\subsection{Mappa concettuale: Modulo 1 $\to$ Modulo 2}

\begin{figure}[H]
\centering
\begin{tcolorbox}[colback=white,colframe=ForestGreen,width=0.92\textwidth,arc=3pt,title=\textbf{Modulo 1}]
Hai implementato e ottimizzato \textbf{una} funzione:\\
\texttt{compute\_partition\_id(key, P)} $\to$ \texttt{partition\_id} $\in [0, P)$\\[3pt]
Varianti: plain (scalare), autovec (GCC), AVX2 (intrinsics)\\
Benchmark: throughput in Mkeys/s, distribuzione tra partizioni
\end{tcolorbox}

\smallskip
$\Big\Downarrow$ \quad \textit{la funzione $h(\mathit{key})$ viene usata qui}

\smallskip
\begin{tcolorbox}[colback=white,colframe=MidnightBlue,width=0.92\textwidth,arc=3pt,title=\textbf{Modulo 2}]
Usi quella funzione \textbf{dentro} un algoritmo completo:
\begin{enumerate}[nosep]
  \item \textbf{Partizionamento}: chiama $h(\mathit{key})$ per ogni record di R e S\\
        $\hookrightarrow$ Histogram + Scatter; chiamata $2\times(N_R+N_S)$ volte
  \item \textbf{Join locale}: non usa più $h(\mathit{key})$; lavora su partizioni con Build+Probe
  \item \textbf{Parallelizzazione}: scegli quali fasi parallelizzare con \texttt{std::thread}
\end{enumerate}
Output: \texttt{join\_count} + \texttt{checksum1} + \texttt{checksum2}\\
Report: strategie + grafici + discussione (max 5 pagine)
\end{tcolorbox}
\end{figure}


% ══════════════════════════════════════════════════════════════
\clearpage
\section{Teoria: Concetti Chiave dalle Lezioni}
% ══════════════════════════════════════════════════════════════

\subsection{Metriche di Performance (Lezione 10)}

\begin{align}
\text{Speedup:}\quad S(p) &= \frac{T_{\text{seq}}}{T_{\text{par}}(p)} &
\text{Efficienza:}\quad E(p) &= \frac{S(p)}{p} = \frac{T_{\text{seq}}}{T_{\text{par}}(p) \cdot p}
\end{align}

\begin{itemize}[nosep]
  \item Speedup ideale: $S(p) = p$ (lineare). In pratica sub-lineare per overhead.
  \item Efficienza ideale: $E(p) = 1$ (100\%).
  \item \textbf{Scalabilità} (relative speedup): $\text{Scalab.}(p) = T_{\text{par}}(1) / T_{\text{par}}(p)$ --- baseline = parallelo con 1 thread.
\end{itemize}

\subsection{Strong vs Weak Scaling (Lezione 10)}

\paragraph{Strong Scaling} (problem size fisso, si variano i processori):
\[
S(p) \leq \frac{1}{f + \frac{1-f}{p}} \quad\text{(Legge di Amdahl)}
\qquad\Rightarrow\qquad
\lim_{p\to\infty} S(p) = \frac{1}{f}
\]
Anche il 5\% di codice seriale limita lo speedup massimo a $20\times$.

\paragraph{Weak Scaling} (problem size cresce con $p$):
\[
S(p) \leq f + p(1-f) \quad\text{(Legge di Gustafson)}
\qquad\qquad
\text{WSE}(p) = \frac{T(PS(1),1)}{T(PS(p),p)}
\]

\subsection{Work-Span Model (Lezione 15)}

Un programma parallelo come DAG: $T_1$ = work totale, $T_\infty$ = span (cammino critico).

\begin{align}
\text{Lower bounds:}\quad T_p &\geq \max\!\left(\frac{T_1}{p},\; T_\infty\right)\\[4pt]
\text{Brent (upper):}\quad T_p &\leq \frac{T_1 - T_\infty}{p} + T_\infty
\end{align}

\textbf{Parallelismo disponibile}: $T_1/T_\infty$ indica il massimo speedup teorico.

\begin{notebox}[Applicazione al progetto]
Ogni fase dell'algoritmo ha un proprio $T_1$ e $T_\infty$. Analizzare il rapporto $T_1/T_\infty$ per ogni fase indica dove il parallelismo è conveniente.
\end{notebox}

\subsection{Tipi di Parallelismo (Lezione 11)}

Per questo progetto si usa \textbf{Data Parallelism}:
\begin{itemize}[nosep]
  \item Pattern \textbf{Map}: partizionamento delle relazioni (ogni thread gestisce un blocco di record)
  \item Pattern \textbf{Reduce}: accumulazione dei risultati parziali
  \item Il join locale per partizione è \emph{embarrassingly parallel}
\end{itemize}

\subsection{Data Partitioning e Workload Balancing (Lezione 14)}

\begin{center}
\begin{tabular}{lll}
\toprule
\textbf{Strategia} & \textbf{Assegnamento} & \textbf{Quando usarla}\\\midrule
Block & $\text{thread}_i \gets [\frac{iN}{p},\;\frac{(i+1)N}{p})$ & Workload regolare (histogram, scatter)\\
Cyclic & $t_i \to \text{thread}_{i \bmod p}$ & Load balancing migliore\\
Block-Cyclic & chunk di dimensione $c$, ciclici & Compromesso\\
Dinamica & coda condivisa / work-stealing & Workload irregolare (join locale)\\
\bottomrule
\end{tabular}
\end{center}

\subsection{C++ Threads (Lezione 13)}

\begin{lstlisting}[style=cpp,caption={Pattern base spawn/join}]
#include <thread>
#include <vector>

std::vector<std::thread> threads;
for (int i = 0; i < nthreads; ++i) {
    threads.emplace_back([i, &shared_data]() {
        // lavoro del thread i
    });
}
for (auto& t : threads) t.join();  // compilare con -pthread
\end{lstlisting}

Regole chiave: ogni thread va joined/detached; \texttt{std::thread} è move-only; evitare oversubscription ($\text{threads} \approx \text{core}$); passare per \texttt{std::ref} con attenzione alla lifetime.

\subsection{Cache e False Sharing (Lezione 5--6)}

\textbf{False Sharing}: thread diversi scrivono su variabili nella stessa cache line (64B) $\Rightarrow$ invalidazione continua.

\begin{lstlisting}[style=cpp,caption={Evitare false sharing con padding}]
// MALE: contatori adiacenti
uint64_t counts[nthreads];  // false sharing!

// BENE: padding per separare le cache lines
struct alignas(64) PaddedCounter { uint64_t value = 0; };
PaddedCounter counts[nthreads];
\end{lstlisting}


% ══════════════════════════════════════════════════════════════
\clearpage
\section{Step 0: Setup dell'Ambiente}
% ══════════════════════════════════════════════════════════════

\subsection{Struttura del progetto}

\begin{lstlisting}[style=plain]
module_2/
  Makefile
  README.md
  include/
    common.hpp          <- tipi, hash, utilita' dal Modulo 1
  src/
    hashjoin_seq.cpp    <- sequenziale (con la TUA hash)
    hashjoin_par.cpp    <- versione parallela
  results/              <- output dei benchmark
  report/
    report.pdf          <- report finale (max 5 pagine)
\end{lstlisting}

\subsection{Makefile}

\begin{lstlisting}[style=make,caption={Makefile}]
CXX      = g++
CXXFLAGS = -O3 -std=c++20 -Wall -Wextra -pthread
INCLUDES = -Iinclude

all: hashjoin_seq hashjoin_par

hashjoin_seq: src/hashjoin_seq.cpp
	$(CXX) $(CXXFLAGS) $(INCLUDES) $< -o $@

hashjoin_par: src/hashjoin_par.cpp
	$(CXX) $(CXXFLAGS) $(INCLUDES) $< -o $@

clean:
	rm -f hashjoin_seq hashjoin_par
\end{lstlisting}

\subsection{Compilazione e test iniziale}

\begin{lstlisting}[style=bash]
g++ -O3 -std=c++20 -Wall hashjoin_seq.cpp -o hashjoin_seq -pthread
./hashjoin_seq -nr 5 -ns 8 -seed 13 -max-key 8 -p 4
./hashjoin_seq -nr 1000000 -ns 2000000 -seed 42 -max-key 100000 -p 64
\end{lstlisting}


% ══════════════════════════════════════════════════════════════
\section{Step 1: Integrare la Funzione di Mapping del Modulo 1}
% ══════════════════════════════════════════════════════════════

Sostituisci \texttt{compute\_partition\_id} nel codice sequenziale:

\begin{lstlisting}[style=cpp,caption={Sostituzione della hash}]
// PRIMA (default del professore):
static inline uint32_t compute_partition_id(uint64_t key, uint32_t p) {
    return (uint32_t)(key & (uint64_t)(p - 1U));
}

// DOPO (la tua dal Modulo 1):
static constexpr uint32_t HASH_A32 = 0x9E3779B9u;

static inline uint32_t compute_partition_id(uint64_t key, uint32_t p) {
    unsigned shift = 32 - __builtin_ctz(p);
    uint32_t k_lo = (uint32_t)(key);
    uint32_t k_hi = (uint32_t)(key >> 32);
    return ((k_lo ^ k_hi) * HASH_A32) >> shift;
}
\end{lstlisting}

\begin{notebox}[Ottimizzazione]
Il valore \texttt{shift} dipende solo da $P$ ed è costante. Conviene pre-calcolarlo nel main e passarlo come parametro per evitare il costo di \texttt{\_\_builtin\_ctz} ad ogni chiamata.
\end{notebox}

La \textbf{stessa identica} funzione deve essere usata nella versione sequenziale e parallela.

% ══════════════════════════════════════════════════════════════
\section{Step 2: Capire il Codice Sequenziale}
% ══════════════════════════════════════════════════════════════

\subsection{Generazione dei dati}
R e S usano seed diverse (\texttt{seed} vs.\ \texttt{seed \^{} 0xdead...}). Le chiavi sono generate con \texttt{splitmix64} e ridotte $\bmod$ \texttt{max\_key}. Più \texttt{max\_key} è piccolo rispetto a $N$, più duplicati ci sono.

\subsection{Histogram --- $O(N)$}
Scansione lineare, scrive su \texttt{hist[pid]} (array di $P$ entry). \textbf{Parallelizzabile} con istogrammi locali + merge.

\subsection{Prefix Sum --- $O(P)$}
Dipendenza loop-carried, ma $P$ è piccolo. Tempo trascurabile $\Rightarrow$ \textbf{non parallelizzare}.

\subsection{Scatter --- $O(N)$}
Scansione lineare con scritture \textbf{random} nell'output. Il cursore \texttt{next[pid]} è condiviso $\Rightarrow$ parallelizzare è il punto più delicato.

\subsection{Join locale (Build + Probe)}
Ogni partizione è \textbf{indipendente}. Il codice usa \texttt{std::unordered\_map} (il professore nota che si può sostituire con strutture più efficienti).

% ══════════════════════════════════════════════════════════════
\section{Step 3: Analisi delle Fasi e Opportunità di Parallelismo}
% ══════════════════════════════════════════════════════════════

\begin{notebox}[Citazione dalla traccia]
``\textit{The goal is not simply to introduce parallelism, but to understand where parallelism is effective and where it is not.}''
\end{notebox}

\begin{center}
\small
\begin{tabular}{llcll}
\toprule
\textbf{Fase} & \textbf{Compl.} & \textbf{Parall.?} & \textbf{Strategia} & \textbf{Note}\\\midrule
Histogram & $O(N)$ & \checkmark~Sì & Istogrammi locali + merge & Facile, buon speedup\\
Prefix Sum & $O(P)$ & $\times$ No & Sequenziale & $P$ piccolo\\
Scatter & $O(N)$ & $\sim$ Possibile & Offset pre-calcolati & Lock-free con local hist\\
Join locale & $O(N)$/part. & \checkmark~Sì & 1 thread per partizione(i) & Embarr.\ parallel\\
Accumulaz. & $O(P)$ & $\times$ No & Sequenziale/reduce & $P$ piccolo\\
\bottomrule
\end{tabular}
\end{center}

Profila il sequenziale per individuare le fasi dominanti (tipicamente scatter + join locale per input grandi).


% ══════════════════════════════════════════════════════════════
\clearpage
\section{Step 4: Implementazione Parallela}
% ══════════════════════════════════════════════════════════════

\subsection{Architettura generale}

\begin{lstlisting}[style=cpp,caption={Struttura del programma parallelo}]
JoinResult partitioned_hash_join_parallel(
    const std::vector<Record>& R, const std::vector<Record>& S,
    uint32_t P, int nthreads)
{
    PartitionedRelation Rpart = partition_relation_parallel(R, P, nthreads);
    PartitionedRelation Spart = partition_relation_parallel(S, P, nthreads);
    return parallel_join(Rpart, Spart, P, nthreads);
}
\end{lstlisting}

\subsection{Histogram parallelo}

Ogni thread calcola un istogramma locale sulla propria porzione (block distribution), poi si fa un merge sequenziale $O(P \times k)$.

\begin{lstlisting}[style=cpp,caption={Histogram parallelo con istogrammi locali}]
// Istogrammi locali (evita false sharing!)
std::vector<std::vector<size_t>> local_hists(nthreads, std::vector<size_t>(P, 0));

std::vector<std::thread> threads;
for (int t = 0; t < nthreads; ++t) {
    threads.emplace_back([&, t]() {
        size_t start = (N * t) / nthreads;
        size_t end   = (N * (t + 1)) / nthreads;
        for (size_t i = start; i < end; ++i)
            ++local_hists[t][compute_partition_id(rel[i].key, P)];
    });
}
for (auto& th : threads) th.join();

// Merge sequenziale
std::vector<size_t> global_hist(P, 0);
for (int t = 0; t < nthreads; ++t)
    for (uint32_t pid = 0; pid < P; ++pid)
        global_hist[pid] += local_hists[t][pid];
\end{lstlisting}

\subsection{Scatter parallelo lock-free (raccomandato)}

Grazie agli istogrammi locali, ogni thread conosce \emph{esattamente} dove scrivere:

\begin{lstlisting}[style=cpp,caption={Offset per-thread-per-partizione + scatter senza lock}]
// Calcolo offset per thread (sequenziale, O(P * nthreads))
std::vector<std::vector<size_t>> thread_offset(nthreads, std::vector<size_t>(P));
for (uint32_t pid = 0; pid < P; ++pid) {
    size_t offset = global_begin[pid];
    for (int t = 0; t < nthreads; ++t) {
        thread_offset[t][pid] = offset;
        offset += local_hists[t][pid];
    }
}

// Scatter parallelo (ZERO lock! regioni di scrittura non overlap)
std::vector<Record> out(N);
for (int t = 0; t < nthreads; ++t) {
    threads.emplace_back([&, t]() {
        size_t start = (N * t) / nthreads;
        size_t end   = (N * (t + 1)) / nthreads;
        auto cursor = thread_offset[t]; // copia locale
        for (size_t i = start; i < end; ++i) {
            uint32_t pid = compute_partition_id(rel[i].key, P);
            out[cursor[pid]++] = rel[i];
        }
    });
}
\end{lstlisting}

\subsection{Join locale parallelo}

Distribuzione ciclica delle partizioni tra i thread (load balancing migliore):

\begin{lstlisting}[style=cpp,caption={Join locale con distribuzione ciclica e padding anti false-sharing}]
struct alignas(64) PaddedResult { JoinResult result{}; };
std::vector<PaddedResult> thread_results(nthreads);

for (int t = 0; t < nthreads; ++t) {
    threads.emplace_back([&, t]() {
        for (uint32_t pid = t; pid < P; pid += nthreads) {
            JoinResult local = join_one_partition(Rpart, Spart, pid);
            thread_results[t].result.join_count += local.join_count;
            thread_results[t].result.checksum1  += local.checksum1;
            thread_results[t].result.checksum2  += local.checksum2;
        }
    });
}
// ... join threads, poi reduce dei risultati
\end{lstlisting}

\subsection{Alternativa: distribuzione dinamica}

\begin{lstlisting}[style=cpp,caption={Scheduler dinamico con atomic counter}]
std::atomic<uint32_t> next_partition{0};
// Ogni thread:
while (true) {
    uint32_t pid = next_partition.fetch_add(1, std::memory_order_relaxed);
    if (pid >= P) break;
    JoinResult local = join_one_partition(Rpart, Spart, pid);
    // accumula...
}
\end{lstlisting}

\subsection{Thread Pool vs Spawn-Join}

Creare thread ad ogni fase introduce overhead. Opzione avanzata: creare i thread una volta e usare \texttt{std::barrier} (C++20) per sincronizzarli tra fasi, riducendo l'overhead $O(p)$ da Amdahl.


% ══════════════════════════════════════════════════════════════
\section{Step 5: Strategia di Validazione (Correctness)}
% ══════════════════════════════════════════════════════════════

\begin{enumerate}
  \item \textbf{Input piccoli} ($N_R, N_S \leq 500$): confronto con \texttt{naive\_join\_verifier} (O($N^2$))
  \item \textbf{Input grandi}: confronto \texttt{join\_count}/\texttt{checksum} tra sequenziale e parallelo
  \item \textbf{Multi-thread}: eseguire con $t = 1, 2, 4, 8, 16$ --- tutti devono dare lo \emph{stesso} risultato
  \item \textbf{Edge cases}: poche partizioni, molti duplicati, thread $>$ partizioni, partizioni vuote
\end{enumerate}

\begin{lstlisting}[style=bash,caption={Script di validazione}]
for T in 1 2 4 8 16; do
    echo "=== Threads=$T ==="
    ./hashjoin_par -nr 1000000 -ns 2000000 -seed 42 \
                   -max-key 1000 -p 64 -t $T
done
# Tutti devono produrre gli STESSI join_count e checksum
\end{lstlisting}

Opzionale ma consigliato: compilare con \texttt{-fsanitize=thread} per esporre data race.


% ══════════════════════════════════════════════════════════════
\section{Step 6: Performance Evaluation}
% ══════════════════════════════════════════════════════════════

Il report richiede \textbf{speedup, strong scaling e weak scaling}.

\subsection{Strong Scaling}
Fissare il problem size, variare i thread. Calcolare:
\[
S(t) = \frac{T_{\text{seq}}}{T_{\text{par}}(t)}, \qquad E(t) = \frac{S(t)}{t}
\]

\subsection{Weak Scaling}
$N_R(t) = t \cdot N_R(1)$, $N_S(t) = t \cdot N_S(1)$. Calcolare:
\[
\text{WSE}(t) = \frac{T(PS(1),1)}{T(PS(t),t)}
\]

\subsection{Breakdown per fase}

\begin{center}\small
\begin{tabular}{l|ccc}
\toprule
\textbf{Fase} & \textbf{1 thread} & \textbf{8 thread} & \textbf{Speedup fase}\\\midrule
Histogram R & xxx ms & xxx ms & x.xx\\
Scatter R & xxx ms & xxx ms & x.xx\\
Histogram S & xxx ms & xxx ms & x.xx\\
Scatter S & xxx ms & xxx ms & x.xx\\
Join locale & xxx ms & xxx ms & x.xx\\
Accumulazione & xxx ms & xxx ms & x.xx\\\midrule
\textbf{TOTALE} & xxx ms & xxx ms & x.xx\\
\bottomrule
\end{tabular}
\end{center}

\subsection{Regole di benchmarking}
\begin{itemize}[nosep]
  \item Almeno 5 ripetizioni, prendere la mediana
  \item Warm-up iniziale
  \item Compilare con \texttt{-O3}
  \item Benchmark su un nodo dell'\texttt{spmcluster}
  \item Non misurare la generazione dei dati
\end{itemize}


% ══════════════════════════════════════════════════════════════
\section{Step 7: Scrivere il Report}
% ══════════════════════════════════════════════════════════════

\begin{center}\small
\begin{tabular}{llc}
\toprule
\textbf{Sezione} & \textbf{Contenuto} & \textbf{Pagine}\\\midrule
1.\ Introduzione & Descrizione del problema e obiettivo & 0.5\\
2.\ Strategia di Parallelizzazione & Analisi per fase, decisioni, come si evitano race/false sharing & 1.5\\
3.\ Risultati Sperimentali & Speedup, weak scaling, breakdown per fase, piattaforma & 2\\
4.\ Discussione & Bottleneck, confronto Amdahl, miglioramenti futuri & 1\\
\bottomrule
\end{tabular}
\end{center}

Grafici da includere: (1) Speedup vs \#threads con linea ideale; (2) Efficienza vs \#threads; (3) Breakdown per fase; (4) Weak scaling efficiency.


% ══════════════════════════════════════════════════════════════
\section{Checklist Finale}
% ══════════════════════════════════════════════════════════════

\begin{itemize}[label=$\square$]
  \item Funzione hash del Modulo 1 integrata in \texttt{compute\_partition\_id}
  \item Versione sequenziale compila e produce output corretto
  \item Versione parallela implementata con \texttt{std::thread}
  \item Correctness: output parallelo $=$ output sequenziale (per tutti i test)
  \item Edge cases testati
  \item Niente race condition (opzionale: testato con \texttt{-fsanitize=thread})
  \item Niente false sharing: padding/allineamento dove necessario
  \item Benchmark strong/weak scaling eseguiti sul cluster
  \item Breakdown per fase misurato e riportato
  \item Makefile presente con istruzioni di compilazione
  \item README con istruzioni per compilare, eseguire, e validare
  \item Report PDF $\leq$ 5 pagine con grafici e discussione
  \item ZIP: \texttt{Modulo2\_NomeCognome.zip} contenente tutto
  \item Email con oggetto ``SPM Modulo2''
\end{itemize}


% ══════════════════════════════════════════════════════════════
\section{Appendice: Comandi Utili \& Riferimenti}
% ══════════════════════════════════════════════════════════════

\subsection{Compilazione}

\begin{lstlisting}[style=bash]
# Standard
g++ -O3 -std=c++20 -Wall -Wextra -pthread src/hashjoin_par.cpp -o hashjoin_par

# Thread sanitizer (debug race condition)
g++ -O1 -g -std=c++20 -fsanitize=thread -pthread src/hashjoin_par.cpp -o hashjoin_par_tsan

# Address sanitizer (debug memory)
g++ -O1 -g -std=c++20 -fsanitize=address -pthread src/hashjoin_par.cpp -o hashjoin_par_asan
\end{lstlisting}

\subsection{Profilazione}

\begin{lstlisting}[style=bash]
# Perf (su Linux/cluster)
perf stat ./hashjoin_par -nr 10000000 -ns 20000000 -seed 42 -max-key 1000000 -p 128 -t 8
perf record ./hashjoin_par ...
perf report
\end{lstlisting}

\subsection{Riferimenti alle lezioni}

\begin{center}\small
\begin{tabular}{ll}
\toprule
\textbf{Argomento} & \textbf{Lezione}\\\midrule
Speedup, Efficienza, Amdahl, Gustafson & \texttt{10-Metrics\_and\_Laws.pdf}\\
Data/Stream/Task Parallelism & \texttt{11-TypesOfParallelism.pdf}\\
C++ auto, move, STL, lambdas & \texttt{12-C++Essentials.pdf}\\
std::thread, mutex, CV, future/promise & \texttt{13-C++ConcurrencyBasics.pdf}\\
Block/Cyclic distribution, Dynamic sched. & \texttt{14-WorkloadBalancing.pdf}\\
Work-Span, Brent's theorem, PRAM & \texttt{15-ModelsOfComputation.pdf}\\
Cache, False Sharing, Roofline model & \texttt{l5 \& 6-Shared-Memory.pdf}\\
SLURM & \texttt{l4-SLURM.pdf}\\
\bottomrule
\end{tabular}
\end{center}

\vfill
\begin{center}
\textit{Buon lavoro! Il professore valuta non solo la correttezza ma la capacità di ragionare sulle scelte di design e di supportarle con dati sperimentali.}
\end{center}

\end{document}
